% !TEX program = lualatex
\documentclass[12pt,a4paper]{article}
\usepackage{amsmath,amssymb,amsthm}
\usepackage{geometry}
\geometry{margin=2.5cm}
\usepackage{hyperref}
\hypersetup{colorlinks=true,linkcolor=blue,citecolor=blue,urlcolor=blue}
\usepackage{booktabs}
\usepackage{lmodern}
\usepackage[T1]{fontenc}
\usepackage{graphicx}

\title{Symmetries Inherent in the Delay Circuit Model}

\author{Noriaki Kihara\\
WF System Co., Ltd.\\
Osaka University, School of Engineering Science (graduated)}

\date{April 2026\quad Research Note\\[4pt]
DOI: \href{https://doi.org/10.5281/zenodo.19534347}{10.5281/zenodo.19534347}}

\begin{document}

\maketitle

\begin{abstract}
% empty
\end{abstract}

\section{Introduction}
% empty

\section{Preliminaries: Summary of Definitions from the Previous Paper}
% empty

\section{Organization of Symmetries}

\subsection{Symmetry of $I$ and $\overline{I}$}

In the previous paper \cite{kihara2026} \S2.2 Property~5, the inversion $\overline{I}$ was defined as
an operation within the same class. That is, $I$ and $\overline{I}$ belong to the same class.

The operation rule of the delay circuit $F_D$ depends solely on the equivalence check between the
input information and the internal information (whether $\mathrm{IN} \neq I$ holds). Replacing $I$
with $\overline{I}$ does not change the equivalence check or the structure of the operation rule.

\paragraph{Definition (Inversion Symmetry).}
The operation rule of the delay circuit $F_D$ is invariant under the transformation $I \to
\overline{I}$. This is called \emph{inversion symmetry}.

This property is a trivial consequence that follows directly from the definition
(\cite{kihara2026} \S2.2) that inversion is an operation within the same class and satisfies
$\overline{\overline{I}} = I$.

\subsection{Definition of the Superclass}

The previous paper \cite{kihara2026} dealt only with operations within the same class (the delay
circuit $F_D$ and the NOT operation circuit). In this section, we define a concept for handling
information belonging to different classes in a unified manner.

A tuple $(A, B, C_1, C_2, \ldots, C_m)$ of information belonging to different classes is defined
as a \textbf{superclass} $\mathcal{I}$. The number of child elements is arbitrary (including zero),
and neither the class nor the value of each child element is constrained. When there are zero child
elements, this corresponds to $\mathrm{null}$.

The superclass is required to satisfy the following properties.

\paragraph{Property 1: Order Independence.}
The superclass does not depend on the order in which child elements appear.
\[
  (A, B) = (B, A)
\]

\paragraph{Property 2: Transparency of null.}
A $\mathrm{null}$ element as a child element of the superclass is equivalent to its absence.
\[
  (A, \mathrm{null}, B) = (A, B)
\]

\paragraph{Property 3: Existence of Inverse Operations.}
For any operation on a superclass $\mathcal{I}$, an inverse operation exists.

The operation satisfying Properties 1 through 3 above is denoted $*$.

\paragraph{Definition of Meta-Information.}
The base class of the superclass has \textbf{meta-information} that is distinct from the
information $I$ itself. Meta-information consists of the following two components:

\begin{itemize}
  \item \textbf{Counter $n$:} Accumulates the number of passes through the delay circuit. Its
        initial value is $\mathrm{null}$.
  \item \textbf{Operation function $*$:} The operation executed when passing through the delay
        circuit. Its initial value is $\mathrm{null}$. The operation function $*$ can reference $n$.
\end{itemize}

Meta-information accompanies information $I$ but is not the value of $I$ itself. Therefore, the
definition of information $I$ in \cite{kihara2026} \S2.2 remains unchanged.

$n$ is an exception to the inverse operation rule: even when an inverse operation is performed, $n$
is preserved.

\subsection{Extended Definition of $F_D$}

The operation rule of $F_D$ defined in \cite{kihara2026} \S2.2 (copying input information to
internal information and passing it to the output) is not changed. When information of the
superclass passes through, $F_D$ performs the following operations on the meta-information of the
passing information, in addition to the original operation rule:

\begin{enumerate}
  \item Increment the counter $n$ by $+1$ (if $n = \mathrm{null}$, set it from $0$ to $1$)
  \item Execute the operation function $*$ if it is not $\mathrm{null}$
\end{enumerate}

The entity that operates on meta-information is $F_D$; the information $I$ itself does not need to
recognize that it ``has passed through $F_D$.''

The NOT operation circuit does not operate on meta-information. This structurally corroborates the
fact stated in \cite{kihara2026} \S2.3 that ``because the NOT operation circuit's operation
involves no delay, it is not counted in the delay count $n$.''

When an operation generates child classes on its own, $n$ can be re-initialized.

\subsection{Consistency between the Superclass and the Delay Circuit}

The superclass $\mathcal{I}$ satisfies the definition of information $I$ in
\cite{kihara2026} \S2.2 (``any information that can be represented numerically''). That is, a
tuple of information $(A, B)$ or $(C_1, C_2, \ldots, C_m)$ belonging to the superclass can each
be treated as a single piece of information $I$.

Therefore, the definitions of the delay circuit $F_D$ and the NOT operation circuit
(\cite{kihara2026} \S2.2--\S2.4) apply to superclass information without modification. The
meta-information operations added in \S3.3 are not a modification of the original operation rule
of $F_D$ but rather the addition of a new operational layer.

\subsection{External Intervention from Outside the System}

The operation rule of $F_D$ defined in the previous paper \cite{kihara2026} and the
meta-information operations in \S3.3 specify the deterministic behavior inside the system. In this
section, we define intervention from outside the system.

\paragraph{External Intervention.}
We allow the operation of forcibly overwriting the internal information of a specific $F_D$ from
outside the system. This operation is performed independently of the operation rules inside the
system.

External intervention is an exception to Property~3 of \S3.2: information that has undergone
external intervention is excluded from the scope of inverse operations. That is, external
intervention introduces irreversibility into the interior of the system.

From the perspective of an observer inside the system, external intervention is an event that
appears suddenly from outside the causal order.

\textit{Note:} This definition suggests that the present model admits a hierarchical structure.
The relationship between an upper-level system and a lower-level system is outside the scope of
this paper.

\section{Discussion}
% empty

\section{Conclusion}

Formulating this model as equations and mapping it to a physical framework are outside the scope
of this paper.

\begin{thebibliography}{9}
\bibitem{kihara2026}
  Noriaki Kihara,
  ``Information-Theoretic Organization of Information Transmission,''
  Research Note, April 2026.
\end{thebibliography}

\end{document}
